\documentclass[UTF8,a4paper,zihao=-4,fontset=fandol]{ctexart}
% Shared formatting for the course report and research manuscript.
\usepackage[margin=2.2cm,headheight=14pt]{geometry}
\usepackage{amsmath,amssymb,bm,booktabs,tabularx,longtable,array}
\usepackage{graphicx,xcolor,tikz}
\usetikzlibrary{arrows.meta,positioning,calc}
\usepackage{caption,float,placeins}
\usepackage[hidelinks,unicode]{hyperref}
\usepackage{url}
\definecolor{signalblue}{HTML}{205A87}
\definecolor{referenceorange}{HTML}{A95822}
\definecolor{softgray}{HTML}{F0F3F5}
\captionsetup{font=small,labelfont=bf,labelsep=quad}
\renewcommand{\arraystretch}{1.2}
\setlength{\tabcolsep}{4.5pt}
\setlength{\parskip}{0.14em}
\setlength{\emergencystretch}{3em}
\linespread{1.22}
\pagestyle{plain}
\setcounter{tocdepth}{2}
\newcommand{\um}{\,\mu\mathrm{m}}
\newcommand{\ns}{\,\mathrm{ns}}
\newcommand{\mV}{\,\mathrm{mV}}
\newcommand{\fC}{\,\mathrm{fC}}
\newcommand{\V}{\,\mathrm{V}}
\newcommand{\sourcepath}[1]{{\footnotesize\nolinkurl{#1}}}
\graphicspath{{figures/}}
\hypersetup{pdfauthor={},pdfsubject={可追溯的阶段研究成果；独立测试与工艺扩展尚待最终验收}}

% Fixed local-ledger snapshot; do not interpret RUNNING as a live process check.
\newcommand{\EvidenceDate}{2026-10-11 05:02\,UTC$-08{:}00$}
\newcommand{\TrainVerified}{64}
\newcommand{\HoldoutVerified}{19}
\newcommand{\HoldoutRunning}{1}
\newcommand{\HoldoutMissing}{12}
\newcommand{\ExtrapVerified}{0}
\newcommand{\ExtrapRunning}{0}
\newcommand{\ExtrapMissing}{32}

\hypersetup{pdftitle={AI辅助Geant4--TCAD联合仿真的DRAM单粒子读出裕量评估}}
\title{\heiti\Large AI辅助Geant4--TCAD联合仿真的\\DRAM单粒子读出裕量评估\\[0.5em]\large 阶段研究稿}
\author{}
\date{\EvidenceDate}

\begin{document}
\maketitle
\begin{abstract}
针对存储节点电荷扰动与读出电路响应相互混合的问题，构建二维访问晶体管与理想电容组成的DRAM探索模型，将Geant4电离沉积总量传递为TCAD等效径迹源，建立匹配中性参考定义的读出裕量及同条件无粒子对照定义的裕量损失。历史D1诊断表明，固定理想参考保留约$12.37\mV$的无粒子读出偏移；对同一无粒子信号状态，匹配参考使差分由$+2.454961\mV$变为$-9.915036\mV$。在独立建立的50$\,\Omega$预充版本中，参考、源隔离及代表点数值检查通过，已完成32个独立粒子事件的64个存储位响应。基于沉积电荷、栅长和角度的四个岭回归模型，在八组物理工况留一交叉验证中，存储0/1的电荷扰动平均折RMSE分别为$0.102888/0.112903\fC$，裕量损失为$0.955160/1.012491\mV$。这些结果支持受约束AI辅助流程、匹配评价方法及条件响应代理在训练工况内的可行性；独立插值留出与尺寸外推尚未完成全量验证，不能据此宣称测试预测精度或真实器件误码概率。模型未作工艺校准，适用范围限于固定粒子能量、所选电磁输运配置和二维等效源近似。
\end{abstract}
\noindent\textbf{关键词：}DRAM；单粒子效应；Geant4；TCAD；匹配参考；岭回归

\section{引言与相关工作}
动态随机存取存储器（DRAM）依靠存储电容的电荷表示信息。访问晶体管导通后，存储节点与位线发生电荷共享，读出系统依据小幅电压差判断存储状态\cite{keeth2001}。粒子入射在硅中产生电子空穴对，随后发生漂移、扩散和复合；生成电荷、端口收集电荷与存储节点净变化具有不同含义\cite{label1996}。因此，粒子使存储电压变化多少，与给定读出过程能否保持正确差分方向，是两个相关但不能等同的问题。

Geant4与TCAD分别适于描述粒子输运和器件电学响应。陈善强、师立勤已将两类工具用于SRAM单粒子效应评估\cite{chen2011}，因此本研究不将工具耦合本身作为首创贡献。机器学习辅助的TCAD单粒子瞬态回归也已有研究\cite{xu2021}。此外，比较重离子、聚焦X射线和脉冲激光诱导响应的研究指出，生成电荷总量相近并不保证器件瞬态一致，空间分布仍影响响应\cite{ryder2021}。这要求联合模型同时说明源项简化、评价对象与验证边界，而非仅给出软件连接和拟合曲线。

本研究以课程提出的AI辅助多特征尺寸存储器评估任务为背景\cite{course}，聚焦三个问题：第一，有限硅区域内的粒子沉积如何转换为可审计的二维TCAD激励；第二，怎样区分无粒子读出偏移与粒子引起的增量；第三，在源项和无粒子基线已知时，低维回归能否描述不同尺寸、粒子和方向下的连续响应。采用简化1T1C单元保留电荷共享与器件输运，暂不模拟完整阵列、真实感放锁存和恢复。

本文贡献是将受约束输入生成、等效源检查、匹配评价和预测冻结组织成一条可追溯的评估链，并以已完成诊断和训练结果检验其中可验证的环节。定量结果分为历史D1诊断与R50主版本两组。前者用于解释参考设计，后者用于多参数响应与训练交叉验证；二者不合并为同一数据集。本文是一份阶段研究稿，独立留出、尺寸外推和氧化层敏感性没有完成的部分，仅说明既定验证方法和当前证据边界。

\section{模型与参数传递}
\subsection{1T1C探索模型及其参数含义}
信号单元由二维平面NMOS、存储电容$C_s$和位线电容$C_b$构成。存储节点记为SN，位线为BL；字线驱动访问晶体管，使SN与BL交换电荷。中性参考采用相同几何、网格、物理模型、电容和时序，其SN初始目标为$0.5\V$。参考器件不施加粒子源。该结构保留了读出驱动对器件和电容的共同作用，但中性状态与信号状态不同，匹配参考并不保证消除所有状态相关响应。

硅横截面宽为$L+0.6\um$、深为$1.2\um$，出平面按$1\um$有效厚度归一化。表\ref{tab:paper-parameters}给出主要设定。四种特征尺寸专指访问晶体管栅长；扫描时名义氧化层、掺杂和集总电容保持不变，因此它们不是四代商业工艺节点。文献支持1T1C和电荷共享原理，而表中的具体数值是继承并冻结的探索设定，尚未由器件电学曲线、版图寄生或实际DRAM数据校准。

\begin{table}[htbp]
\centering\small
\caption{R50主模型的主要参数；取值为项目设定，未作工艺校准}
\label{tab:paper-parameters}
\begin{tabularx}{\textwidth}{@{}p{3.2cm}X@{}}
\toprule
参数 & 取值或定义\\
\midrule
栅长$L$、名义栅氧厚度 & $0.18/0.25/0.35/0.45\um$；$5\,\mathrm{nm}$\\
体区、源漏掺杂 & 硼$10^{17}\,\mathrm{cm}^{-3}$；砷高斯峰值$10^{20}\,\mathrm{cm}^{-3}$\\
源漏深度参数 & $0.12\um$处指定$10^{17}\,\mathrm{cm}^{-3}$，横向因子0.4；不等于净掺杂结深\\
栅功函数、温度 & $4.6\,\mathrm{eV}$；$300\,\mathrm{K}$\\
存储、位线电容 & $C_s=10\,\mathrm{fF}$；$C_b=100\,\mathrm{fF}$\\
信号、中性初态 & 信号$0.1/0.9\V$对应存储0/1；中性$0.5\V$\\
预充与字线 & BL目标$0.5\V$；R50预充通路$50\,\Omega$；WL终值$1.8\V$\\
\bottomrule
\end{tabularx}
\end{table}

器件求解耦合泊松方程、电子与空穴连续性方程及外部电路，采用Fermi统计、OldSlotboom带隙变窄、掺杂相关迁移率、高场饱和、法向电场迁移率修正，以及SRH和Auger复合。瞬态采用后向欧拉离散。理想钳位在$0.6\ns$释放，$0.95\ns$核验初态；粒子在$1\ns$入射，$2.8\ns$提取保持阶段指标。BL在$2.8$--$2.9\ns$预充，$3\ns$将预充电阻切至$10^{20}\,\Omega$；WL在$3.05$--$3.10\ns$上升，$4\ns$采样。受控初态使尺寸和位值的比较有共同起点，但不构成真实写入过程的验证。

\subsection{输运几何与等效径迹}
Geant4 10.7.3采用option4电磁构造器和步长限制\cite{g4em}。粒子为$5\,\mathrm{MeV}$的alpha或proton，方向角为$0^\circ/30^\circ$；材料仅含硅和外部真空。硅横截面与TCAD对应，入射点为$(L+0.45,0,0)\um$，方向为$(-\sin\theta,\cos\theta,0)$，$+y$指向硅体内部。有限直径迹长度为
\begin{equation}
\ell=\begin{cases}
1.2\um,&\theta=0^\circ,\\
\min\!\left(\dfrac{1.2}{\cos\theta},\dfrac{L+0.45}{\sin\theta}\right)\um,&\theta=30^\circ.
\end{cases}
\label{eq:paper-length}
\end{equation}
式中$L$以$\mu\mathrm{m}$计。例如$L=0.18\um$、$30^\circ$时从侧边退出，$\ell=1.26\um$；该直径迹不等于散射后的实际输运路程。固定表面入口也意味着本文比较的是约定入射位置，不能推广到整个单元的空间敏感区分布。

\subsection{沉积量到TCAD源项的转换与检查}
设第$j$步总沉积和非电离沉积分别为$E_{{\rm dep},j}$与$E_{{\rm NIEL},j}$，定义
\begin{align}
E_{\rm ion}&=\sum_j\max(0,E_{{\rm dep},j}-E_{{\rm NIEL},j}),\\
Q_{\rm dep}&=q\frac{10^6E_{\rm ion}[\mathrm{MeV}]}{W[\mathrm{eV}]},\qquad
W=3.6\,\mathrm{eV},\quad q=1.602176634\times10^{-19}\,\mathrm{C}.
\label{eq:paper-qdep}
\end{align}
$W$是硅中生成一对电子空穴的平均能量约定\cite{ryder2021}。$Q_{\rm dep}$表示一类载流子的生成电荷尺度，不是产生的净电荷，也不是SN最终损失的电荷。

将总量沿式\eqref{eq:paper-length}的直径迹均匀分配，线电荷密度为$\lambda_{\rm eq}=Q_{\rm dep}/\ell$，横向高斯尺度为$0.03\um$。程序字段\texttt{LET}采用$\mathrm{pC}/\mu\mathrm{m}$，与常用质量阻止本领单位不同。逐步沉积记录用于溯源；实际激励是总量等效源，未忠实传递三维逐步沉积及其时空分布。

对TCAD生成率$G_{\rm ion}$作时空积分，并将二维归一化厚度计入体积元：
\begin{equation}
Q_{\rm src}=q\int\!\!\int_{\Omega}G_{\rm ion}(\mathbf r,t)\,\mathrm dV\,\mathrm dt,
\qquad \frac{|Q_{\rm src}-Q_{\rm dep}|}{Q_{\rm dep}}\leq5\%.
\label{eq:paper-source}
\end{equation}
相对门槛适用于正沉积事件，合法零沉积单独检查近零源项。侧边高斯截断不能通过改小目标沉积量来通过检查。该积分检验转换总量；对等效分布的物理适用性，仍须保留降维近似的限制。

\section{AI辅助流程及评价指标}
\subsection{受约束交互与证据链}
流程如图\ref{fig:paper-workflow}所示。AI辅助编写和审查参数生成、结果提取与证据核验程序；自然语言入口接受约定中文模板，映射到粒子、方向、尺寸和存储位等结构化参数，再检查范围与单位。R50演示支持预览、审计和定位已有证据，实际求解由冻结输入和固定控制器执行。该结果证明受约束请求的参数映射，不能写成任意自然语言生成任意模型或自主开启新求解。

\begin{figure}[htbp]
\centering
% Conceptual workflow, not a measured result or screenshot.
\resizebox{\linewidth}{!}{%
\begin{tikzpicture}[
 block/.style={draw=signalblue,rounded corners=2pt,fill=softgray,
 text width=4.2cm,minimum height=1.5cm,align=center,font=\small},
 arrow/.style={-{Latex[length=2.4mm]},thick,draw=signalblue},
 node distance=0.7cm and 0.6cm]
\node[block] (request) {限定中文请求\\尺寸、粒子、角度\\独立事件数};
\node[block,right=of request] (matrix) {受约束的工况清单\\单位、范围、版本与种子};
\node[block,right=of matrix] (transport) {Geant4 硅内电离沉积\\$E_{\rm ion}\ \rightarrow\ Q_{\rm dep}$\\有限径迹及单位检查};
\node[block,below=of transport] (tcad) {TCAD 器件与电路求解\\等效径迹与受控初态\\读出瞬态与端口检查};
\node[block,left=of tcad] (metric) {匹配中性参考与同位对照\\$D,\ M,\ \Delta Q,\ \Delta M$};
\node[block,left=of metric] (model) {按工况分组的响应拟合\\冻结预测后检验未见响应};
\draw[arrow] (request)--(matrix);
\draw[arrow] (matrix)--(transport);
\draw[arrow] (transport)--(tcad);
\draw[arrow] (tcad)--(metric);
\draw[arrow] (metric)--(model);
\node[below=0.40cm of metric,align=center,font=\small,text width=14.5cm]
 {每一层绑定输入版本、运行记录和来源哈希；科学质量门槛与执行状态分别记录。};
\end{tikzpicture}}

\caption{AI辅助评估流程及证据传递。自然语言映射、粒子事件、TCAD响应和冻结预测具有不同产物；本文只将已归档并通过相应检查的数据作为结果。}
\label{fig:paper-workflow}
\end{figure}

每个响应关联模型版本、粒子事件和输入标识。检查覆盖初态、预充、源项积分、端口电流闭合以及所需匹配关系；求解结束标志不足以证明物理评价可用。确定性重跑保留原失败身份并关联替代归档，不增加独立样本。流程将AI建议、固定规则与求解器结果分开，便于定位错误来源；本文未设人工操作对照组，因此不报告AI带来的效率提升比例。

\subsection{粒子增量与读出参考的分工}
分别用$e$、$c$、$n$表示粒子事件、同条件同存储位无粒子对照和中性参考。在保持阶段$t_h=2.8\ns$，定义
\begin{equation}
\Delta Q=C_s\big[V_{{\rm SN},e}(t_h)-V_{{\rm SN},c}(t_h)\big].
\label{eq:paper-dq}
\end{equation}
$\Delta Q$以fC报告，负值表示存储电容相对对照减少电荷。两种位值对同一方向扰动的敏感性不同，因此不能把全部$\Delta Q$称为正的电荷损失。

对位值$b\in\{0,1\}$，在$t_r=4\ns$定义
\begin{align}
D_e&=V_{{\rm BL},e}(t_r)-V_{{\rm BL},n}(t_r),\\
M_e&=(2b-1)D_e,\qquad
M_c=(2b-1)\big[V_{{\rm BL},c}(t_r)-V_{{\rm BL},n}(t_r)\big],\\
\Delta M&=M_c-M_e.
\label{eq:paper-margin}
\end{align}
$D$、$M$、$\Delta M$以mV报告，$\Delta M>0$表示粒子引起裕量下降。中性参考定义读出差分，位匹配对照提取粒子增量；二者不可互换。同一匹配参考在$\Delta M$中相消，也不意味着可以省略它，因为绝对裕量和判据分类仍取决于参考。

主判据阈值为$\epsilon=5\mV$：$M>\epsilon$为正确方向，$M<-\epsilon$为判据误读，$|M|\leq\epsilon$为临界；另按$2.5/10\mV$检查阈值敏感性。阈值为项目约定，未标定实际感放失调和噪声。数值检查不足以支持边界分类的响应保留“数值不确定”，不能并入正确或误读。本文的误读是固定采样时刻的模型判据，不等同于已锁存的持久位翻转。

\subsection{矩阵设计与独立事件单位}
设计包含4种栅长、2种粒子、2个角度，每工况4个独立事件，共64事件。每个事件分别计算存储0和1，共128个配对响应。训练采用$L=0.18/0.35\um$，含32事件、64响应；$L=0.25\um$为尺寸插值留出，$L=0.45\um$为尺寸外推，各16事件、32响应。训练与测试按尺寸分离，配对位值及同事件的氧化层或数值变体均不增加粒子事件分母。该设计检验特定尺寸变化，不构成对任意能量、入射位置或材料体系的泛化验证。

\section{参考验证与多参数响应}
\subsection{历史D1诊断：固定参考中的无粒子偏移}
D1在$L=0.25\um$下完成9项有效无粒子工况，比较存储初态、字线驱动和参考方式。中性SN的实际读前电压为$0.500055125\V$，二电容封闭共享预测差分仅$0.001544\mV$，TCAD的$4\ns$差分却为$12.369997\mV$。理想共享关系
\begin{equation}
V_f=\frac{C_sV_s(t_0)+C_bV_b(t_0)}{C_s+C_b}
\label{eq:paper-sharing}
\end{equation}
假定两电容间封闭交换且最终等势；实际访问器件和字线驱动使这一假设并不充分。中性工况外接电容净增电荷为$1.360530\fC$，除以$C_s+C_b$对应约$12.368453\mV$，与上述差额一致。按该诊断的符号约定，栅、体端积分分别为$+1.464107$和$-0.103577\fC$，与净增电荷闭合。栅端项是位移电流，不能解释为载流子穿越栅氧；端口闭合也不能唯一拆分重叠电容、沟道形成和耗尽区变化的贡献。

关闭字线时偏移接近零；将上升时间由50 ps增至200 ps，已报告精度内的终点偏移未消失。这支持读出驱动参与偏移形成，而非将偏移归为新产生的粒子电荷。对无粒子构造的同一信号状态，固定理想参考给出$+2.454961\mV$，匹配中性参考给出$-9.915036\mV$，见图\ref{fig:paper-d1}。物理双支路与离线相减满足$0.1\mV$预设容差，说明两种参考实现一致；它不是独立实测校准。该无粒子状态仅匹配历史事件的读前SN电压，不能声称复现其完整载流子历史。所有D1数值保留历史条件，不作为R50偏移量或R50粒子响应。

\begin{figure}[htbp]
\centering
\includegraphics[width=0.88\textwidth]{d1-reference-comparison.pdf}
\caption{历史D1中同一无粒子信号状态的参考比较。参考改变读出差分的符号，说明绝对判据依赖参考定义；该比较不是粒子诱发持久位翻转的证明，也不代替R50基线。}
\label{fig:paper-d1}
\end{figure}

\subsection{R50参考、预充与数值检查}
原100$\,\Omega$版本出现$13.297\,\mu\mathrm V$的预充误差，超过$10\,\mu\mathrm V$门槛。半步长与细网格诊断没有使其通过，故保留失败记录。R50将预充通路改为$50\,\Omega$，建立独立输入版本并重跑主线。对同一固定事件的独立探针比较，BL预充误差由$13.297362$变为$4.731204\,\mu\mathrm V$，见图\ref{fig:paper-precharge}；这项比较不能回写旧结果，也不能用放宽门槛解释通过。

R50已完成20项参考计算，并通过源隔离和预登记代表点的数值检查。图\ref{fig:paper-reference}展示四尺寸的无粒子预击前初态及代表尺寸的读出过程。另由参考验收记录分别核对物理双支路与离线方法的信号侧、参考侧和差分，最大电压差约$6\times10^{-15}\V$；该数值反映同模型计算一致性，不代表实际测量精度。源积分以原Geant4沉积总量为目标，代表点细化采用$0.5\mV$的读出量容差及$0.01\fC$的电荷差容差。通过支持已检查输入的实现与离散稳定性，不代表整个参数域均完成收敛验证，也不证明模型已符合真实DRAM工艺。

\begin{figure}[htbp]
\centering
\includegraphics[width=0.88\textwidth]{r50-precharge.pdf}
\caption{同一固定事件的预充电阻比较。柱高为$3\ns$时$|V_{\rm BL}-0.5\V|$：原100$\,\Omega$失败项为$13.297362\,\mu\mathrm V$，50$\,\Omega$独立探针为$4.731204\,\mu\mathrm V$，门槛保持$10\,\mu\mathrm V$。图中不含参考节点曲线，也不改写原失败身份。}
\label{fig:paper-precharge}
\end{figure}
\begin{figure}[htbp]
\centering
\includegraphics[width=0.88\textwidth]{r50-reference.pdf}
\caption{R50无粒子基准。左图为四个栅长的存储0、中性和存储1共12项预击前SN电压；右图为$L=0.18\um$三种状态在$3$--$4\ns$的BL曲线。数据来自已验收参考归档；图中没有粒子响应，也不是物理双支路与离线方法的差值图。}
\label{fig:paper-reference}
\end{figure}

\subsection{已完成训练矩阵的连续响应与分类}
R50训练部分的32个独立事件、64个位响应已完成。图\ref{fig:paper-training-response}按训练工况展示连续量，位值配对用于观察同一沉积对两种状态的影响。按冻结$5\mV$主判据，64响应中48项为正确方向，10项为判据误读，4项临界，2项数值不确定。连续量可用于拟合，并不意味着每个分类已确定；2项不确定不能强行转为正确或误读，也不能因不利于结论而删除。

\begin{figure}[htbp]
\centering
\includegraphics[width=\textwidth]{r50-training-response.pdf}
\caption{R50训练连续响应与Geant4沉积电荷$Q_{\rm dep}$的关系。两行分别为$2.8\ns$的$\Delta Q$和$4\ns$的$\Delta M$，两列对应存储0/1；包括$L=0.18/0.35\um$、alpha/proton及$0^\circ/30^\circ$。每工况4个独立事件，两种位值合计64响应，测试尺寸未混入。}
\label{fig:paper-training-response}
\end{figure}

这些计数表明，同一评估链能够保留裕量充足、临界和方向相反的响应，而非只输出成功求解的二元标签。当前每一尺寸、粒子、角度及位值层仅4个独立事件，适于方法与响应探索，不足以估计低概率误码或截面。尤其不能把64个位响应当作64次独立入射，或把判据误读比例解释为真实存储器辐射失效率。

独立留出与外推按既定分母继续验收，表\ref{tab:paper-progress}记录本文快照的完成范围。运行中和未完成响应不进入最终预测误差；训练结果和预先冻结的测试预测，也不能补作缺失测试真值。
\begin{table}[htbp]
\centering\small
\caption{固定证据快照中的矩阵完成范围；状态描述不代表独立测试评分完成}
\label{tab:paper-progress}
\begin{tabular}{llrrrr}
\toprule
阶段 & 栅长 / $\mu$m & 计划响应 & 已验证 & 在途 & 待提交 \\
\midrule
训练 & 0.18 / 0.35 & 64 & 64 & 0 & 0 \\
域内留出 & 0.25 & 32 & 19 & 1 & 12 \\
域外外推 & 0.45 & 32 & 0 & 0 & 32 \\
\bottomrule
\end{tabular}

\end{table}
\FloatBarrier

\section{拟合外推及讨论}
\subsection{条件响应代理与分组验证}
预测目标为式\eqref{eq:paper-dq}的$\Delta Q$和式\eqref{eq:paper-margin}的$\Delta M$。对每个位值与目标分别建立模型，共四个回归器。特征为
\begin{equation}
\mathbf x=\left(Q_{\rm dep},\;Q_{\rm dep}/L,\;Q_{\rm dep}\sin\theta\right),
\label{eq:paper-features}
\end{equation}
依次采用fC、$\mathrm{fC}/\mu\mathrm m$、fC为单位；模型使用训练折均值和总体标准差标准化。$Q_{\rm dep}/L$引入尺寸相关尺度，角度项提供简单方向依赖，但这些组合是低维代理特征，不构成对微观收集机制的唯一解释。

岭回归最小化
\begin{equation}
\sum_{i\in\mathcal T}\left[y_i-\beta_0-\boldsymbol\beta^{\mathsf T}\widetilde{\mathbf x}_i\right]^2
+\lambda\|\boldsymbol\beta\|_2^2,
\label{eq:paper-ridge}
\end{equation}
其中截距不受惩罚。候选$\lambda$为$0$、$0.01$、$0.1$、$1$、$10$。训练数据按$(L,\text{粒子},\theta)$分成八个物理工况，逐组留出；每组四次重复及两种位值保持同组，避免同事件跨训练与验证。选择指标是八折RMSE的算术平均，四个模型均选得$\lambda=0.1$。标准化只使用该折训练数据，最终拟合才使用全部训练响应。比较基准是对应训练折目标均值。

\subsection{训练交叉验证结果及其解释}
表\ref{tab:paper-training-cv}和图\ref{fig:paper-training-cv}给出训练阶段结果。存储0/1的$\Delta Q$平均折RMSE为$0.102888/0.112903\fC$，训练均值基准为$2.279967/2.515710\fC$；$\Delta M$为$0.955160/1.012491\mV$，均值基准为$21.458870/23.383861\mV$。在当前工况中，引入沉积、尺寸和角度特征比常数预测保留了更多响应变化信息。

\begin{table}[htbp]
\centering\small
\caption{八组物理工况留一交叉验证的平均折RMSE；属于训练阶段模型选择，不是独立测试误差}
\label{tab:paper-training-cv}
\begin{tabular}{clrr}
\toprule
存储位 & 目标（单位） & 岭回归 RMSE & 均值基准 RMSE \\
\midrule
0 & $\Delta Q$ (fC) & 0.102888 & 2.279967 \\
0 & $\Delta M$ (mV) & 0.955160 & 21.458870 \\
1 & $\Delta Q$ (fC) & 0.112903 & 2.515710 \\
1 & $\Delta M$ (mV) & 1.012491 & 23.383861 \\
\bottomrule
\end{tabular}

\end{table}
\begin{figure}[htbp]
\centering
\includegraphics[width=0.9\textwidth]{r50-training-cv.pdf}
\caption{训练交叉验证中岭回归与训练均值基准的比较。不同目标分别保留fC和mV单位；误差来自八个工况折的RMSE平均，不能当作全部留出点合并后的RMSE或独立外推精度。}
\label{fig:paper-training-cv}
\end{figure}

上述误差是用于选择正则化参数的同一交叉验证过程的结果，不能视为无偏的最终泛化估计。特别是八折留出一个物理工况，并不等同于完整留出一个从未参与训练的栅长。误差低于均值基准支持该参数化表示的训练内可行性；它尚未回答$0.25\um$插值和$0.45\um$外推是否可靠，也未验证预测分类能否稳定越过$5\mV$边界。

\subsection{独立预测的冻结条件与待验收范围}
模型、标准化参数以及64个测试位响应的预测已在独立测试响应提交前冻结。最终评分必须复用这些预测，分别给出插值留出与尺寸外推的MAE、RMSE、最大绝对误差及分层残差，不按测试表现重新调参，也不筛去不利工况。设计未规定必须优于均值基准才纳入结果；真实完成且通过对应审计的测试应按预定规则报告，包括预测失败或优势不明显的情况。

本代理的输入包含真实Geant4沉积，目标的定义和绝对裕量还依赖已知无粒子基线，因此它是给定源项和基线条件下的TCAD响应代理。它没有从任意粒子描述直接预测器件结果，也没有替代所有Geant4与无粒子TCAD计算。本文快照尚未完成独立测试全量验收，不给出其最终误差、预测对真值散点图或外推通过结论。

\subsection{模型适用边界与后续接入}
尺寸变化可能同时改变输运、端口耦合和电荷共享动态，本模型只改变栅长并保持多项参数固定，无法分离真实工艺节点缩放中的全部因素。当前质子结果限于所选电磁配置下的电离响应，不能据此讨论核反应次级粒子或核反冲主导的单粒子效应。硅/真空几何、等效直径迹和二维归一化也未覆盖完整材料栈及真实三维敏感体积。源积分守恒和数值一致性是必要检查，但不能替代实验或可信工艺模型的校准。

已规划的$4/6\,\mathrm{nm}$氧化层比较采用固定硅入口事件，只能作为TCAD参数敏感性分析；它不构成跨完整工艺模型的粒子输运重算，也不属于当前栅长预测的氧化层外推验证。其余扩展、边界数值检查及R50原生图形取证尚待各自完成归档，不在本文画占位结果。后续结果应保持R50版本、事件身份和匹配条件一致，并保留失败及确定性重跑的对应关系，避免把不同模型或重复求解混成额外物理证据。

\section{结论}
本文建立了从Geant4沉积总量、TCAD等效源到DRAM读出裕量的可追溯评估链。历史D1结果表明，无粒子读出过程可产生约$12.37\mV$的固定参考偏移；同一信号状态在匹配参考下发生差分符号变化。由此明确，中性参考负责定义读出裕量，同位无粒子对照负责识别粒子增量，存储电荷扰动与最终判据应分别报告。

独立R50版本已完成参考、源隔离、代表点数值检查及32事件的64个训练响应。基于三个沉积与几何特征的四个岭回归器，在八组物理工况留一交叉验证中得到约$0.10$--$0.11\fC$的电荷扰动平均折RMSE及$0.96$--$1.01\mV$的裕量损失平均折RMSE，均低于对应训练均值基准。该结果支持当前条件响应代理的训练内可行性；独立插值与尺寸外推仍需按冻结预测完成验证。

AI的已证实作用是辅助构建受约束输入、提取与审计流程，未证明任意自然语言自主仿真或操作效率优势。研究结论限于未校准1T1C探索模型、固定$5\,\mathrm{MeV}$粒子条件和当前电磁输运及二维近似，不推导真实器件误码概率、核反冲效应或商业节点可靠性。后续工作应以完成既定独立测试、边界检查和参数敏感性归档为依据扩充结论，避免用计划或训练拟合代替验证。


\FloatBarrier
\bibliographystyle{unsrt}
\bibliography{references}

\appendix
\section{证据来源与复核范围}
正文由固定证据快照生成，证据日期为\EvidenceDate。D1、原100$\,\Omega$版本与R50分别保留模型身份；路径、运行标识、输入和原始输出哈希以及图表字段对应关系由随稿证据包记录\cite{project}。表\ref{tab:paper-evidence}列出主要读取对象。训练交叉验证图表来自已冻结模型，不能用实时测试进度替换其中的误差定义。

\begin{table}[htbp]
\centering\small
\caption{主要证据对象及其用途}
\label{tab:paper-evidence}
\begin{tabularx}{\textwidth}{@{}p{2.3cm}X p{4.8cm}@{}}
\toprule
对象 & 归档入口 & 本文使用范围\\
\midrule
D1诊断 & \sourcepath{experiments/dram-readout-offset-d1/data/summary.csv} & 历史无粒子偏移与参考比较\\
R50冻结模型 & \sourcepath{experiments/dram-see-v3/r50-main/data/predictions/models.json} & 特征、分组、正则化及训练交叉验证\\
R50预测归档 & \sourcepath{experiments/dram-see-v3/r50-main/data/predictions/} & 预测先于独立测试的冻结证据\\
本文结果快照 & \sourcepath{tables/snapshot.tex}及随稿图表来源索引 & 状态分母、完成范围和版本对应\\
\bottomrule
\end{tabularx}
\end{table}

对文稿的复核包括来源一致性、公式和单位、样本分母、阶段界限及PDF排版。本文整理不启动新的粒子输运或器件求解，也不将待完成扩展作为结果。模型与预测的科学结论仍以正文所列已完成检查为限。
\end{document}
